\chapter{Futures sur indices boursiers et sur actions dans le monde}\label{ch:m1:index-and-single-stock-futures-worldwide}

En 2025, les bourses du monde ont échangé 119 milliards de contrats à terme
et d'options. Soixante-trois sur cent ont changé de mains en Asie-Pacifique,
vingt en Amérique du Nord, moins de quatre en Europe~; et quatre-vingts sur cent
portaient sur des actions. Un lecteur qui n'a vu que les contrats de Chicago du
\cref{ch:m1:futures-contracts-and-exchanges} n'a vu qu'un cinquième du tableau,
mesuré en nombre de contrats. Mesuré en dollars, le tableau est très différent,
et cet écart est la première leçon du chapitre. La seconde tient à ce que le
future sur actions a des cousins~: des contrats sur les dividendes qu'un indice
versera, des contrats sur son rendement total, des contrats sur des actions
isolées. Chacun isole l'un des termes que la formule de juste valeur de la
\cref{prop:m1:basis-roll-and-delivery:fv} regroupe.

\section{Compter les contrats, compter l'argent}

\begin{definition}[Point d'indice, contrats mini et micro]\label{def:m1:index-and-single-stock-futures-worldwide:point}
Un \emph{point d'indice}\index{point d'indice} est une unité de l'indice coté~; le
multiplicateur d'un contrat donne à ce point sa valeur en argent. Un \emph{contrat
mini}\index{contrat mini} ou \emph{contrat micro}\index{contrat micro}
est une version d'un contrat à terme existant dont le multiplicateur est réduit, typiquement
au cinquième ou au dixième, puis encore au dixième, de sorte que le même marché puisse être
négocié pour des tailles plus petites~; les minis finissent souvent par être bien plus
liquides que les contrats dont ils sont issus.
\end{definition}

\begin{omfigure}
\begin{tikzpicture}
\begin{axis}[width=10.5cm,height=4.8cm,xbar,bar width=9pt,
  xlabel={contrats échangés en 2025 (milliards)},
  symbolic y coords={Other,Europe,Latin America,North America,Asia-Pacific},ytick={Other,Europe,Latin America,North America,Asia-Pacific},yticklabels={{Autres},{Europe},{Amérique latine},{Amérique du Nord},{Asie-Pacifique}},
  tick label style={font=\small},label style={font=\small},
  xmin=0,xmax=92,xmajorgrids,grid style={gray!25},enlarge y limits=0.12,
  nodes near coords,nodes near coords style={font=\footnotesize,fill=white,inner sep=1pt,/pgf/number format/.cd,fixed,fixed zerofill,precision=1},
  table/col sep=comma]
\addplot[fill=omDef!60,draw=omDef] table[y=region,x=contracts_bn]{figdata/markets-1/22-index-and-single-stock-futures-worldwide/regions.csv};
\end{axis}
\end{tikzpicture}

\omcaption{Contrats à terme et options échangés en bourse en 2025, par région, en
nombre de contrats~: 119.3 milliards au total, 42~\% de moins qu'en 2024. Un contrat
n'est pas une unité de risque~: l'un d'eux peut valoir 300\,000~\$ et un autre
quelques centaines de dollars. Données~: statistiques annuelles de l'association
professionnelle du secteur.}\label{fig:m1:index-and-single-stock-futures-worldwide:regions}
\end{omfigure}

\begin{dated}{2026-09}{Ce que disent les statistiques de volume, et ce qu'elles ne disent pas}\label{dat:m1:index-and-single-stock-futures-worldwide:fia}
L'association professionnelle des contrats à terme a compté 119.29 milliards de contrats en 2025~:
88.65 milliards d'options, en baisse de 50.3~\% sur l'année, et 30.64 milliards de contrats à terme, en hausse de
8.6~\%. Les contrats sur actions s'élevaient à 94.79 milliards, en baisse de 48.8~\%, ce qui explique
l'essentiel de la chute. Par région~: Asie-Pacifique 75.59 milliards ($-55.6\%$), Amérique du Nord
24.50 ($+24.0\%$), Amérique latine 11.39, Europe 4.38, autres 3.43. Les comptages
pèsent chaque contrat de la même façon. Une division par deux du total mondial en une année où
le volume des contrats à terme a augmenté signifie que le nombre est dominé par des options sur
actions de très petite taille sur quelques marchés, et par les décisions de leurs régulateurs sur
la taille des contrats et les échéances, et non par le montant de risque transféré.
\end{dated}

\section{Trois références régionales}

\begin{center}
\small
\begin{tabular}{llll}
\toprule
 & E-mini S\&P 500 & Future EURO STOXX 50 & Future Hang Seng Index \\
\midrule
Bourse & Chicago & Francfort & Hong Kong \\
Multiplicateur & 50~\$ & 10~\euro & 50~HK\$ \\
Tick & 0.25 (12.50~\$) & 1 (10~\euro) & 1 (50~HK\$) \\
Échéances & trimestrielles & trimestrielles & mensuelles, trimestrielles \\
Dernier jour & 3\up{e}~vendredi, ouverture & 3\up{e}~vendredi, midi & avant-dernier jour ouvré \\
Règlement final & cours d'ouverture & moyenne sur dix minutes & moyenne des cotations sur 5~min \\
\bottomrule
\end{tabular}
\end{center}

\begin{dated}{2026-09}{Horaires}\label{dat:m1:index-and-single-stock-futures-worldwide:hours}
Le contrat à terme EURO STOXX 50 se négocie de 02 h 10 à 22 h 00, heure de Francfort. Le
contrat à terme Hang Seng se négocie de 09 h 15 à 12 h 00 et de 13 h 00 à 16 h 30,
heure de Hong Kong, puis dans une séance hors cotation de 17 h 00 à 03 h 00, dont
les transactions relèvent de la compensation du jour suivant. Ses mois d'échéance
sont le mois au comptant, les trois mois civils suivants et les trois mois trimestriels
suivants, avec des contrats de juin et de décembre à plus longue échéance.
\end{dated}

Chacun des trois est l'instrument de couverture des produits dérivés sur actions
de sa région et le premier prix à réagir à l'annonce d'une nouvelle lorsque le
marché au comptant est fermé. Ensemble, ils sont ouverts presque en permanence, du matin
asiatique au soir américain (\cref{fig:m1:index-and-single-stock-futures-worldwide:relay}), et une société de trading qui en
négocie un en continu doit savoir ce qui est ouvert ailleurs~: la liquidité de chaque
contrat augmente nettement lorsque son propre marché au comptant est ouvert, et de
nouveau lorsque celui d'une autre région ouvre.

\begin{omfigure}
\begin{tikzpicture}
\begin{axis}[width=12cm,height=3.8cm,
  xlabel={heure (UTC), mercredi 16 septembre 2026},
  ytick={1,2,3},yticklabels={ES,FESX,HSI},
  tick label style={font=\small},label style={font=\small},yticklabel style={font=\small\ttfamily},
  xmin=0,xmax=24,ymin=0.4,ymax=3.6,xtick={0,3,6,9,12,15,18,21,24},xmajorgrids,grid style={gray!25},
  unbounded coords=jump,table/col sep=comma]
\addplot[omDef,line width=7pt] table[x=hour,y=row]{figdata/markets-1/22-index-and-single-stock-futures-worldwide/relay.csv};
\end{axis}
\end{tikzpicture}

\omcaption{Le relais sur une journée, en UTC. Hong Kong fait une pause pour le déjeuner~;
rien n'est ouvert entre 21 h 00 et 22 h 00. Les barres européennes et américaines se
décalent d'une heure par rapport à UTC, et par rapport à Hong Kong, lorsque
l'heure d'été prend fin. Les horaires d'ES relèvent d'une hypothèse des données de test du
chapitre, et non d'une cotation de la règle de la bourse. Données~: la construction du
chapitre.}\label{fig:m1:index-and-single-stock-futures-worldwide:relay}
\end{omfigure}

\section{Contrats à terme sur dividendes}

\begin{definition}[Contrat à terme sur dividendes]\label{def:m1:index-and-single-stock-futures-worldwide:divfut}
Un \emph{contrat à terme sur dividendes}\index{contrat à terme sur dividendes} sur un indice se règle à la
somme des dividendes bruts ordinaires versés par les composantes de l'indice au cours
d'une période de contrat, en général un an, exprimée en \omterm{def:m1:index-and-single-stock-futures-worldwide:point}{points d'indice}. Avant l'échéance
son prix est le prix de marché de ce flux de dividendes.
\end{definition}

\begin{dated}{2026-09}{Le contrat à terme sur dividendes de l'indice EURO STOXX 50}\label{dat:m1:index-and-single-stock-futures-worldwide:fexd}
Valeur du contrat 100~\euro par point de dividende de l'indice~; pas de cotation de 0.1 point, soit 10~\euro.
Cotation~: les cinq prochains mois trimestriels, les deux mois semestriels suivants
et les huit prochains mois de décembre. La négociation prend fin à 12 h 00 CET le
troisième vendredi du mois d'échéance~; le prix de règlement-livraison final est le
total cumulé des dividendes bruts pertinents versés par les composantes de l'indice
pour la période du contrat, tel que défini par le fournisseur de l'indice.
\end{dated}

Le contrat existe en raison d'un sous-produit. Les banques qui vendent des produits structurés
sur un indice se retrouvent longues des dividendes futurs de l'indice pour des années
à venir, un risque qu'elles ne peuvent couvrir avec des contrats à terme sur indices, dont l'exposition aux dividendes
elles détiennent déjà. La vente de \omterm{def:m1:index-and-single-stock-futures-worldwide:divfut}{contrats à terme sur dividendes} transfère ce risque. Les acheteurs sont
des fonds qui jugent le prix bas~: une série de \omterm{def:m1:index-and-single-stock-futures-worldwide:divfut}{contrats à terme sur dividendes} s'est souvent négociée
sous toute prévision raisonnable des dividendes pour les années éloignées, la
décote étant le prix que les banques paient pour se débarrasser du risque. La pente de la série
n'est donc pas une prévision (\cref{fig:m1:index-and-single-stock-futures-worldwide:strip}).

\begin{proposition}[Dividendes implicites]\label{prop:m1:index-and-single-stock-futures-worldwide:implied}
À partir d'un contrat à terme sur indice $F$, de temps $T$ restant avant l'échéance, du cours au comptant $S$ et du taux de financement
$r$, les dividendes que le marché valorise entre maintenant et l'échéance sont
\[
D^{\text{impl}} = S(1+rT) - F
\]
(réinvestissement négligé). Un contrat à terme sur indice acheté revient à être vendeur de ces dividendes~: si
les dividendes attendus baissent de $\delta$ points, la fair value du contrat à terme monte de
$\delta$.
\end{proposition}

\begin{omfigure}
\begin{tikzpicture}
\begin{axis}[width=10cm,height=4.6cm,ybar,bar width=14pt,
  xlabel={année du contrat},ylabel={contrat à terme sur dividendes (points)},
  xtick=data,x tick label style={/pgf/number format/1000 sep=},
  tick label style={font=\small},label style={font=\small},
  ymin=150,ymax=178,ymajorgrids,grid style={gray!25},
  nodes near coords,nodes near coords style={font=\footnotesize,fill=white,inner sep=1pt,/pgf/number format/.cd,fixed,fixed zerofill,precision=1},
  table/col sep=comma]
\addplot[fill=omDef!60,draw=omDef] table[x=year,y=points]{figdata/markets-1/22-index-and-single-stock-futures-worldwide/strip.csv};
\end{axis}
\end{tikzpicture}

\omcaption{Une série illustrative de contrats à terme annuels sur dividendes. La croissance implicite est
de $+2.1\%$ pour la première année, puis négative, soit $-0.5\%$ par an pendant les
cinq années~: la forme que produit l'offre de couverture, quelles que soient les prévisions des analystes.
L'axe vertical ne part pas de zéro. Données~: le
tutoriel.}\label{fig:m1:index-and-single-stock-futures-worldwide:strip}
\end{omfigure}

\section{Contrats à terme de rendement total}

\begin{definition}[Contrat à terme de rendement total]\label{def:m1:index-and-single-stock-futures-worldwide:trf}
Un \emph{contrat à terme de rendement total}\index{contrat à terme de rendement total} est un contrat à terme coté
qui réplique un \omterm{def:m1:financing:trs}{total return swap} sur actions
(\cref{def:m1:delta-one-instruments:swap})~: le porteur reçoit le rendement de cours de l'indice
et ses distributions, et paie le taux de référence overnight majoré
d'un spread. Il est coté et négocié en termes de ce spread, en points de base par an~;
la bourse convertit ce spread en prix de contrat à terme en \omterm{def:m1:index-and-single-stock-futures-worldwide:point}{points d'indice} à l'aide
du cours de clôture de l'indice, des distributions courues et du financement couru.
\end{definition}

\begin{dated}{2026-09}{Le contrat à terme de rendement total sur l'indice EURO STOXX 50}\label{dat:m1:index-and-single-stock-futures-worldwide:tesx}
Multiplicateur 10~\euro par \omterm{def:m1:index-and-single-stock-futures-worldwide:point}{point d'indice}. Coté en spread annualisé, en points de base
avec une décimale, variation minimale de 0.5 point de base. Les sous-jacents sont
l'indice de prix, l'indice de distributions du fournisseur et le taux à court terme en euros
sans écart, comme taux de financement. Les distributions et le financement courent depuis le
lancement et s'ajoutent au prix~; leurs variations quotidiennes passent par la \omterm{def:m1:clearing-and-settlement:margin}{marge de
variation}. Les échéances s'étendent sur près de dix ans~: les 21 prochains mois trimestriels
et cinq décembres supplémentaires.
\end{dated}

Dans le contrat à terme ordinaire sur indice, financement et dividendes sont inclus dans un même nombre,
la base, et le porteur est exposé aux deux. Le \omterm{def:m1:index-and-single-stock-futures-worldwide:trf}{contrat à terme de rendement total} verse
les dividendes réalisés et prélève un taux overnight réalisé, de sorte que ce qui reste à négocier est
le spread seul~: le prix à terme du financement en actions, la
même grandeur que la richesse du roll dans
la \cref{def:m1:basis-roll-and-delivery:roll}, fixée pour des années au lieu d'un
trimestre et compensée. Les banques l'utilisent pour sortir de leur bilan le financement de leurs
couvertures en actions~; les fonds l'utilisent pour verrouiller le coût de financement de
une position de long terme.

\section{Contrats à terme sur actions}

\begin{definition}[Contrat à terme sur action]\label{def:m1:index-and-single-stock-futures-worldwide:ssf}
Un \emph{contrat à terme sur action}\index{contrat à terme sur action} est un contrat à terme portant sur une
action, \omterm{def:m1:basis-roll-and-delivery:delivery}{réglé en espèces} ou par livraison des actions. C'est un \omterm{def:m1:delta-one-instruments:deltaone}{delta-one}
enveloppe (\cref{ch:m1:delta-one-instruments})~: financement inclus dans le prix, aucun
emprunt de titres nécessaire pour être short, et les opérations sur titres gérées par
l'ajustement du contrat opéré par la bourse.
\end{definition}

Leur sort est national. Ils se négocient activement en Inde et sur les
bourses européennes de produits dérivés, où une bourse cote des contrats à terme et des options
sur plus de 1\,200 sous-jacents issus de plus de vingt pays, la majeure partie du
volume étant constitué de transactions motivées par les dividendes et le financement, entre professionnels
autour des dates de dividende. Aux États-Unis, la seule bourse qui leur était consacrée
a cessé de négocier en septembre 2020 et a retiré son agrément, les combinaisons
de swaps et d'options ayant accaparé l'activité. Lorsqu'une taxe sur les
transactions s'applique aux actions et non aux contrats à terme, ou lorsque la vente
à découvert d'actions est restreinte, le \omterm{def:m1:index-and-single-stock-futures-worldwide:ssf}{contrat à terme sur action} est l'instrument naturel~; lorsque
aucune des deux ne s'applique, c'est le swap qui l'emporte.

\section{Tutoriel~: dividendes implicites et strip}\label{tut:m1:index-and-single-stock-futures-worldwide:implied}

\begin{tutorial}
\textbf{Objectif.} Extraire les dividendes implicites d'un contrat à terme sur indice, lire la
croissance d'un strip de dividendes, évaluer un \omterm{def:m1:index-and-single-stock-futures-worldwide:trf}{contrat à terme de rendement total} à partir de son
spread, et trouver quels contrats sont ouverts à un instant donné. \textbf{État final~:} les
trois figures de données de ce chapitre.
\begin{tutsteps}
\item \textbf{Dividendes implicites et croissance.}
\omcode{code/markets-1/22-index-and-single-stock-futures-worldwide/python/div_implied.py}{4}{11}{Dividendes implicites d'un contrat à terme, et croissance implicite d'un strip.}\label{lst:m1:index-and-single-stock-futures-worldwide:implied}
\item \textbf{Un \omterm{def:m1:index-and-single-stock-futures-worldwide:trf}{contrat à terme de rendement total} à partir de son spread.}
\omcode{code/markets-1/22-index-and-single-stock-futures-worldwide/python/div_implied.py}{20}{24}{Du spread négocié à un prix en points d'indice.}\label{lst:m1:index-and-single-stock-futures-worldwide:trf}
\item \textbf{Sessions en UTC.} Heures locales, fuseaux horaires, sessions qui franchissent
minuit, semaines qui commencent le dimanche.
\omcode{code/firm/calendar/firm_calendar.py}{29}{37}{L'intervalle UTC d'une session qui s'ouvre à une date locale donnée.}\label{lst:m1:index-and-single-stock-futures-worldwide:session}
\end{tutsteps}
\textbf{Pour aller plus loin.} Ajouter les jours fériés de chaque bourse pour l'année
en cours et recalculer les heures de l'année pendant laquelle les trois contrats sont
ouverts ensemble.
\end{tutorial}

\section{Construction~: le calendrier mondial des contrats}\label{bld:m1:index-and-single-stock-futures-worldwide:calendar}

\begin{build}
\bfield{Rôle} Les stratégies de la société miniature qui s'étendent sur plusieurs régions
doivent savoir, pour tout horodatage, quels contrats sont ouverts, et le simulateur doit
refuser les ordres sur les contrats fermés.
\bfield{Interface} \texttt{load(path)} renvoyant les sessions~;
\texttt{Session.interval(local\_date)}~; \texttt{is\_open(sessions, root,
when)}~; \texttt{open\_roots(sessions, when)}~; \texttt{utc\_intervals(sessions,
root, day)}.
\bfield{Règles} Les sessions sont stockées en heure locale de la bourse, avec un fuseau
horaire IANA~; toutes les réponses sont en UTC, et le changement d'heure provient de la base
des fuseaux horaires, jamais d'un décalage saisi à la main. Une session dont la clôture ne
se situe pas après l'ouverture se termine le jour local suivant. La semaine est définie par les
jours de semaine où une session \emph{s'ouvre}. Les jours fériés sont hors champ et documentés comme tels.
\bfield{Données} \texttt{data/markets-1/sessions\_sample.csv}~: cinq lignes~; la
ligne ES est une hypothèse, et le fichier l'indique.
\bfield{Tests d'acceptation} \texttt{code/firm/calendar/tests/}~: l'ouverture européenne qui
se décale d'une heure en UTC entre septembre et décembre~; la pause déjeuner de Hong Kong et sa
session à cheval sur minuit~; les bornes du vendredi et du dimanche~; l'ensemble des contrats
ouverts à trois instants~; le découpage à un jour UTC.
\bfield{Pour aller plus loin} Jours fériés et demi-journées par bourse~; clôtures anticipées
le jour d'échéance, à partir du référentiel des contrats.
\end{build}

\begin{omsources}
\item FIA, \emph{ETD Volume --- December 2025} (totaux annuels par catégorie et
région).
\item Eurex, pages produits \emph{EURO STOXX 50 Index Futures}, \emph{EURO
STOXX 50 Index Dividend Futures (FEXD)}, \emph{EURO STOXX 50 Index Total
Return Futures (TESX)} et la circulaire sur la transition du taux de financement du TESX vers
\euro STR flat~; \emph{Single Stock Futures}.
\item HKEX, \emph{Hang Seng Index Futures}, résumé du contrat.
\item US Securities and Exchange Commission, ordonnance faisant droit à la demande de retrait
d'agrément de OneChicago LLC, \emph{Federal Register}, 18 février
2021.
\end{omsources}

\section{Exercices}

\begin{exercise}[$\star$]\label{exo:m1:index-and-single-stock-futures-worldwide:1}
Donner le notionnel d'un contrat de chacun des trois indices de référence aux niveaux
de 6\,000, 5\,400 et 26\,000, dans sa devise propre, et la valeur d'un mouvement
de 1~\%.
\end{exercise}

\begin{exercise}[$\star$]\label{exo:m1:index-and-single-stock-futures-worldwide:2}
L'indice vaut 5\,400, le contrat à terme à trois mois 5\,412.0 et le taux de financement
2.5~\%. Donner les dividendes implicites sur le trimestre.
\end{exercise}

\begin{exercise}[$\star$]\label{exo:m1:index-and-single-stock-futures-worldwide:3}
Un fonds achète 50 \omterm{def:m1:index-and-single-stock-futures-worldwide:divfut}{contrats à terme sur dividendes} de décembre 2027 à 171.5. Les
dividendes de l'année s'élèvent à 158.2 points. Donner le P\&L.
\end{exercise}

\begin{exercise}[$\star\star$]\label{exo:m1:index-and-single-stock-futures-worldwide:4}
À partir du strip de la \cref{fig:m1:index-and-single-stock-futures-worldwide:strip}
(168.0, 171.5, 169.0, 166.0, 164.5, 163.5), calculer la croissance implicite de chaque année
et la croissance annualisée sur les cinq années. Un analyste prévoit $+4\%$ par
an. Le strip est-il «~faux~»~?
\end{exercise}

\begin{exercise}[$\star\star$]\label{exo:m1:index-and-single-stock-futures-worldwide:5}
À partir de la \cref{fig:m1:index-and-single-stock-futures-worldwide:relay}~: pendant combien
d'heures de ce jour UTC les trois contrats sont-ils ouverts ensemble, et pendant quelle
heure aucun n'est ouvert~? Quelles barres évoluent en décembre, et dans quel sens~?
\end{exercise}

\begin{exercise}[$\star\star$]\label{exo:m1:index-and-single-stock-futures-worldwide:6}
Un fonds est acheteur de 200 \omterm{def:m1:index-and-single-stock-futures-worldwide:trf}{contrats à terme de rendement total} à deux ans de l'échéance, achetés
à un spread de 40 points de base, l'indice étant à 5\,400. Le spread passe à 55.
Donner le P\&L et expliquer son signe.
\end{exercise}

\begin{exercise}[$\star\star\star$]\label{exo:m1:index-and-single-stock-futures-worldwide:7}
\emph{Programmation.} Avec \texttt{firm\_calendar}, trouver, pour le mercredi 16 décembre
2026, l'heure UTC à laquelle aucun des trois contrats n'est ouvert, et la comparer
à celle de septembre. Puis lister chaque instant de la semaine du 14 septembre où
l'ensemble des contrats ouverts change.
\end{exercise}

\begin{exercise}[$\star\star\star$]\label{exo:m1:index-and-single-stock-futures-worldwide:8}
\emph{Trouver la faille.} «~L'Asie-Pacifique traite 63~\% des dérivés mondiaux,
donc c'est là que 63~\% du risque mondial est transféré, et là qu'une nouvelle société de
trading doit s'implanter en premier~». Utiliser
l'\cref{dat:m1:index-and-single-stock-futures-worldwide:fia}.
\end{exercise}

\section{Problème~: La courbe des dividendes}

\begin{problem}[{Problème du week-end --- négocier ce que l'indice va verser}]\label{pb:m1:index-and-single-stock-futures-worldwide:1}
L'indice est à 5\,400. Les contrats à terme annuels sur dividendes pour 2026 à 2031 se négocient à
168.0, 171.5, 169.0, 166.0, 164.5 et 163.5 \omterm{def:m1:index-and-single-stock-futures-worldwide:point}{points d'indice} (100~\euro par point). Le
contrat à terme sur indice a un multiplicateur de 10~\euro.

\medskip\noindent\textbf{Partie I --- Lecture du strip.}
\begin{enumerate}
\item Donner le rendement de dividende implicite pour 2027.
\item Donner la croissance implicite de 2026 à 2027.
\item Donner la croissance implicite annualisée de 2026 à 2031.
\item Les dividendes ont progressé de 4~\% par an pendant une décennie. Donner le dividende 2031 suivant
cette tendance, ainsi que la décote du contrat à terme 2031 par rapport à celui-ci.
\item Qui vend le contrat à terme 2031 avec cette décote, et pourquoi~?
\end{enumerate}

\medskip\noindent\textbf{Partie II --- Une position.}
Un fonds achète 200 contrats à terme 2031 à 163.5.
\begin{enumerate}[resume]
\item Donner le notionnel en euros.
\item Donner le P\&L si les dividendes 2031 se matérialisent selon la tendance de 4~\%.
\item Donner le P\&L s'ils sont amputés de 30~\% par rapport au niveau de 2026, comme
dans une forte récession.
\item Le contrat est réévalué quotidiennement. En période de crise, il tombe à 110 en un mois
alors que l'échéance 2031 n'est qu'à cinq ans. Donner la \omterm{def:m1:clearing-and-settlement:margin}{marge de variation}, et expliquer pourquoi un
investisseur de long terme ne peut pas l'ignorer.
\end{enumerate}

\medskip\noindent\textbf{Partie III --- Les dividendes dans le contrat à terme sur indice.}
Un desk est acheteur de 500 contrats à terme sur indice échéant en décembre 2027 à titre de couverture.
\begin{enumerate}[resume]
\item Les dividendes 2027 attendus sont réduits de 20~\%. Donner la variation en \omterm{def:m1:index-and-single-stock-futures-worldwide:point}{points d'indice}.
\item Donner le P\&L du desk résultant de cette seule nouvelle valorisation, avec son signe.
\item Combien de \omterm{def:m1:index-and-single-stock-futures-worldwide:divfut}{contrats à terme sur dividendes} 2027 neutraliseraient cette exposition~?
\item Expliquer pourquoi un desk de produits structurés finit par être acheteur de dividendes.
\item Pourquoi ce même desk préférerait-il un \omterm{def:m1:index-and-single-stock-futures-worldwide:trf}{contrat à terme de rendement total} au contrat à terme
ordinaire pour sa couverture~?
\end{enumerate}

\medskip\noindent\textbf{Partie IV --- Jugement.}
\begin{enumerate}[resume]
\item Le strip est-il une prévision~? Qu'est-ce que c'est~?
\item Les \omterm{def:m1:index-and-single-stock-futures-worldwide:divfut}{contrats à terme sur dividendes} ont chuté bien davantage que les dividendes n'ont finalement été réduits
lors des crises passées. Qu'est-ce que cela révèle de ceux qui les détiennent~?
\item Qu'est-ce qui joue aujourd'hui le rôle de «~sous-jacent~» pour le contrat à terme 2031,
et quelles implications cela a-t-il pour sa couverture~?
\item Doit-on s'attendre à ce que les \omterm{def:m1:index-and-single-stock-futures-worldwide:divfut}{contrats à terme sur dividendes} d'une seule action soient plus
ou moins liquides que le contrat sur indice~? Pourquoi existent-ils~?
\item Énoncer le \emph{résultat nommé}~: la croissance implicite des dividendes, sur la première année
et annualisée sur cinq ans.
\item En une phrase~: que dissocie un \omterm{def:m1:index-and-single-stock-futures-worldwide:divfut}{contrat à terme sur dividendes}~?
\end{enumerate}
\end{problem}

\section{Questions d'entretien}

\begin{interviewq}[$\star$ \iqroles{trader, chercheur}]\label{iq:m1:index-and-single-stock-futures-worldwide:1}
Pouvez-vous citer le principal contrat à terme sur indice actions de trois régions, et une
façon dont leurs échéances diffèrent~?
\end{interviewq}

\begin{interviewq}[$\star$ \iqroles{trader, chercheur, banque}]\label{iq:m1:index-and-single-stock-futures-worldwide:2}
Qu'est-ce qu'un \omterm{def:m1:index-and-single-stock-futures-worldwide:divfut}{contrat à terme sur dividendes}, et qui l'utilise~?
\end{interviewq}

\begin{interviewq}[$\star\star$ \iqroles{trader, banque}]\label{iq:m1:index-and-single-stock-futures-worldwide:3}
Vous êtes acheteur d'un contrat à terme sur indice. Une grande valeur de l'indice annule son
dividende. Qu'advient-il de votre P\&L~?
\end{interviewq}

\begin{interviewq}[$\star\star$ \iqroles{trader, banque, chercheur}]\label{iq:m1:index-and-single-stock-futures-worldwide:4}
Qu'est-ce qu'un \omterm{def:m1:index-and-single-stock-futures-worldwide:trf}{contrat à terme de rendement total} vous permet de négocier
que ne permet pas un contrat à terme sur indice ordinaire~?
\end{interviewq}

\begin{interviewq}[$\star\star$ \iqroles{développeur, chercheur}]\label{iq:m1:index-and-single-stock-futures-worldwide:5}
Comment stockeriez-vous les heures de négociation de deux cents contrats sur vingt
bourses~?
\end{interviewq}

\begin{interviewq}[$\star\star\star$ \iqroles{chercheur, trader}]\label{iq:m1:index-and-single-stock-futures-worldwide:6}
Un tableau de classement range les bourses selon le nombre de contrats négociés. Comment
le transformeriez-vous en un outil qu'une société de trading peut utiliser~?
\end{interviewq}
